\chapter{Resultados Esperados (a priori)}
\label{chap:resultados_esperados}

Si el plan se cumple como está escrito en introducción, problema y objetivos, al cerrar la tesis deberían aparecer resultados tangibles en cuatro frentes.

\begin{enumerate}
    \item \textbf{Código entregable:} librería o configuración TypeScript/JavaScript que reemplace los plugins Capacitor prioritarios por respuestas parametrizables, más ejemplos en el proyecto Ionic de referencia.
    \item \textbf{Relato del método de cobertura:} evidencia de qué escenarios se corrieron, qué ramas nuevas se vieron en el reporte y cómo se comparó contra la línea base descrita en metodología.
    \item \textbf{MCP auxiliar funcionando:} servidor pequeño con herramientas declaradas (tests + cobertura + Sonar opcional) y notas de uso acorde a \cite{mcp-spec2025}.
    \item \textbf{Números compilados:} tablas o gráficos con cobertura antes/después, tiempo invertido en preparar mocks y, si existía Sonar/ESLint en el entorno, lecturas complementarias obtenidas con comandos reproducibles.
    \item \textbf{Límites escritos:} lista explícita de plugins cubiertos y advertencia de que todo corresponde a unitarias en Node, no a integración sobre hardware.
\end{enumerate}

La hipótesis se juzga contra ese paquete de evidencias. Los porcentajes exactos se informarán recién cuando existan mediciones reales; no los adelanto en el plan.
